\documentclass[11pt,a4paper]{article}

% Shared front-matter data. Edit here, not in the report bodies.

\newcommand{\coursename}{Blockchain Technology (BITS F452)}
\newcommand{\groupnumber}{6}
\newcommand{\projectname}{DVote}

% Columns are sized to their content so the rules always span the full table.
\newcommand{\groupmemberstable}{%
  \small
  \begin{tabular}{@{}lll@{}}
    \toprule
    \textbf{ID No.} & \textbf{Name} & \textbf{Email} \\
    \midrule
    2024A7PS0148H & Varad Kailas Thombre & \texttt{f20240148@hyderabad.bits-pilani.ac.in} \\
    2024A7PS0121H & Naman Singhal        & \texttt{f20240121@hyderabad.bits-pilani.ac.in} \\
    2024ADPS0190H & Shivam Goyal         & \texttt{f20240190@hyderabad.bits-pilani.ac.in} \\
    2024A7PS0101H & Rishit Agarwal       & \texttt{f20240101@hyderabad.bits-pilani.ac.in} \\
    2024A7PS0113H & Vihaan Kashyap       & \texttt{f20240113@hyderabad.bits-pilani.ac.in} \\
    2024A7PS0041H & Palak Agarwal        & \texttt{f20240041@hyderabad.bits-pilani.ac.in} \\
    2023B4A70727H & Chaitanya Keyal      & \texttt{f20230727@hyderabad.bits-pilani.ac.in} \\
    2024B4A70830H & Wamika Sehgal        & \texttt{f20240830@hyderabad.bits-pilani.ac.in} \\
    \bottomrule
  \end{tabular}%
}


\newcommand{\reporttitle}{\projectname: A Decentralized Voting System}
\newcommand{\reportsubtitle}{Evaluation 1: Design Report}
\newcommand{\submissiondate}{September 2026}

% Shared preamble for both evaluation reports.
% Compile with: latexmk -pdf eval1.tex   (see Makefile)

\usepackage[T1]{fontenc}
\usepackage[utf8]{inputenc}
\usepackage{lmodern}
% The .docx template uses Arial throughout; Helvetica is its metric-compatible
% LaTeX equivalent.
\usepackage{helvet}
\renewcommand{\familydefault}{\sfdefault}
% Margins and header spacing match the template: 0.5in all round, with room for
% the institute logo that sits in the header of every page.
\usepackage[a4paper,margin=0.5in,headheight=88pt,headsep=14pt,includehead]{geometry}
\usepackage{microtype}

\usepackage{amsmath,amssymb}
\usepackage{graphicx}
\usepackage{booktabs}
\usepackage{longtable}
\usepackage{array}
\usepackage{ragged2e}
\usepackage{enumitem}
\usepackage{xcolor}
\usepackage{caption}
\usepackage{fancyhdr}
\usepackage{titlesec}
\usepackage{listings}
\usepackage{csquotes}
\usepackage{tikz}
\usetikzlibrary{arrows.meta,positioning,shapes.geometric,fit,backgrounds,calc}
% Graphs are drawn from the committed CSVs in bench/results, never by hand.
\usepackage{pgfplots}
\usepackage{pgfplotstable}
\pgfplotsset{compat=1.18}
\usepgfplotslibrary{groupplots}

\usepackage[backend=biber,style=numeric,sorting=none]{biblatex}
\addbibresource{references.bib}

\usepackage[hidelinks,breaklinks]{hyperref}

% Keep the PDF free of metadata: no title, author, creator or producer
% strings, no timestamps, no TeX version banner and no random trailer ID.
\hypersetup{pdftitle={}, pdfauthor={}, pdfsubject={}, pdfkeywords={},
            pdfcreator={}, pdfproducer={}}
\pdfinfoomitdate=1
\pdftrailerid{}
\pdfsuppressptexinfo=-1

\graphicspath{{figures/}}

% ---------------------------------------------------------------------------
% Draft placeholders
%
% \stub{...} marks unwritten content. It renders as a visible red box, so an
% unfinished section is obvious on the page. Delete the \stub call as you write
% the prose that replaces it; when none are left, the report is done.
% ---------------------------------------------------------------------------
\newcommand{\stub}[1]{%
  \par\noindent
  \begingroup
  \color{red!65!black}\small\itshape
  \fbox{\parbox{\dimexpr\linewidth-2\fboxsep-2\fboxrule}{TODO: #1}}%
  \endgroup
  \par\medskip
}

% Inline marker for a single value still to be filled in.
\newcommand{\note}[1]{\textcolor{blue!55!black}{[#1]}}

% ---------------------------------------------------------------------------
% Layout
% ---------------------------------------------------------------------------
% Every page carries the institute logo in the header, as the template does.
\newcommand{\institutelogo}{\includegraphics[width=5.81in]{bits-logo}}

\pagestyle{fancy}
\fancyhf{}
\fancyhead[C]{\institutelogo}
\fancyfoot[C]{\thepage}
\renewcommand{\headrulewidth}{0pt}

% Title page: logo header, no page number.
\fancypagestyle{titlestyle}{%
  \fancyhf{}%
  \fancyhead[C]{\institutelogo}%
  \renewcommand{\headrulewidth}{0pt}%
}

% The plain style is used on the first page of the document and by \chapter-like
% breaks; give it the same header so no page is missing the logo.
\fancypagestyle{plain}{%
  \fancyhf{}%
  \fancyhead[C]{\institutelogo}%
  \fancyfoot[C]{\thepage}%
  \renewcommand{\headrulewidth}{0pt}%
}

\titleformat{\section}{\large\bfseries}{\thesection.}{0.6em}{}
\titleformat{\subsection}{\normalsize\bfseries}{\thesubsection}{0.6em}{}

\setlist[itemize]{topsep=2pt,itemsep=2pt,parsep=0pt,leftmargin=1.4em}
\setlist[enumerate]{topsep=2pt,itemsep=2pt,parsep=0pt,leftmargin=1.6em}

\captionsetup{font=small,labelfont=bf}

% A page holding only floats is centred vertically by default, which strands a
% figure in the middle of an otherwise empty page. Push those pages to the top.
\makeatletter
\setlength{\@fptop}{0pt}
\makeatother

% Fixed-width column types for tables that need wrapping.
\newcolumntype{L}[1]{>{\RaggedRight\arraybackslash}p{#1}}
\newcolumntype{R}[1]{>{\RaggedLeft\arraybackslash}p{#1}}

\lstset{
  basicstyle=\ttfamily\footnotesize,
  breaklines=true,
  frame=single,
  framesep=4pt,
  rulecolor=\color{black!30},
  numbers=left,
  numberstyle=\tiny\color{black!50},
  keywordstyle=\bfseries,
  showstringspaces=false,
  captionpos=b,
}

% Solidity is not a built-in listings language; define enough to be readable.
\lstdefinelanguage{Solidity}{
  keywords={pragma,solidity,contract,library,interface,function,modifier,event,
      emit,struct,enum,mapping,memory,storage,calldata,public,private,
      internal,external,view,pure,payable,returns,return,require,revert,
      if,else,for,while,new,constructor,immutable,constant,is,using,
      address,uint256,bytes32,bool,string,error},
  keywordstyle=\bfseries,
  comment=[l]{//},
  morecomment=[s]{/*}{*/},
  commentstyle=\color{black!55}\itshape,
  string=[b]",
  stringstyle=\color{black!70},
  morestring=[b]',
}

% ---------------------------------------------------------------------------
% Title page
% ---------------------------------------------------------------------------
\newcommand{\maketitlepage}{%
  % The title page needs no hyperref anchor, and suppressing it avoids a
  % duplicate page.1 destination once front matter switches to roman numerals.
  \hypersetup{pageanchor=false}%
  \begin{titlepage}
    \thispagestyle{titlestyle}
    \centering
    \vspace*{1.5cm}

    {\Huge\bfseries \reporttitle\par}
    \vspace{0.8cm}
    {\Large \reportsubtitle\par}

    \vspace{1.5cm}
    {\large \coursename\par}
    \vspace{0.3cm}
    {\large \textbf{Group \groupnumber}\par}

    \vspace{1.2cm}
    {\large\bfseries Group Members\par}
    \vspace{0.4cm}
    \groupmemberstable

    \vfill
  \end{titlepage}
  % Front matter in roman numerals so the title page and the body do not both
  % claim page 1 (which produces duplicate hyperref anchors).
  \pagenumbering{roman}%
  \hypersetup{pageanchor=true}%
}

% Ends the front matter and starts the numbered body.
\newcommand{\startbody}{\clearpage\pagenumbering{arabic}}


\begin{document}

\maketitlepage

\tableofcontents
\startbody

% ===========================================================================
\section{Introduction}
% Source material: README.md, docs/architecture.md
% Keep this short, one page. State what the system is, who it is for, and
% what the reader will find in the rest of the report.
% ===========================================================================

\stub{One paragraph on what \projectname{} is: a decentralized application that
  replaces a centrally administered election with a smart-contract-governed one.}

\stub{One paragraph on why this domain was chosen over the other options in the
  problem statement, and what specifically makes voting a good fit for a
  blockchain rather than a database.}

\stub{One short paragraph mapping the rest of the report: Section 2 analyses the
  domain, Section 3 gathers stakeholder requirements, Section 4 proposes the
  design, Section 5 defines how the project will be evaluated, and Section 6 sets
  out the implementation plan.}

% ===========================================================================
\section{Domain and Problem Analysis}
% Source material: README.md ("The problem"), docs/architecture.md section 2
% This section must establish that a centralized solution is genuinely
% inadequate; otherwise the whole project is unjustified.
% ===========================================================================

\subsection{The centralized status quo}

\stub{Describe how elections are administered today (paper ballots, EVMs,
  centrally hosted e-voting). Be concrete and fair: state what these systems do
  well before criticising them.}

\subsection{Trust concentrated in a single operator}

A conventional electronic voting system asks the electorate to trust one operator
with four separable responsibilities at once: the voter roll, the ballot box, the
count, and the audit log. Table~\ref{tab:trust} states each assumption and how a
blockchain-based design removes the need to make it.

\begin{table}[htbp]
  \centering
  \caption{Trust assumptions in a centralized system, and how \projectname{} removes each.}
  \label{tab:trust}
  \begin{tabular}{@{}L{6.0cm}L{9.7cm}@{}}
    \toprule
    \textbf{Assumption the voter must make} & \textbf{How \projectname{} removes it}                                                                      \\
    \midrule
    The voter roll is correct
                                            & The roll is committed as a Merkle root on-chain before voting opens and cannot be changed afterwards.       \\
    Nobody voted twice
                                            & Each ballot carries a zero-knowledge nullifier; the contract rejects duplicates without learning who voted. \\
    My ballot was counted
                                            & Every ballot is an on-chain transaction the voter can locate for themselves.                                \\
    The count is honest
                                            & Tallies are incremented by contract code, not by an operator.                                               \\
    The audit log was not edited
                                            & The log is the chain.                                                                                       \\
    \bottomrule
  \end{tabular}
\end{table}

\stub{Expand each row into prose. For each one, name a concrete failure mode of
  the centralized system: what exactly can a dishonest or compromised operator
  do, and why can a voter not detect it?}

\subsection{Why paper is not the answer either}

\stub{Paper ballots remove some of these problems and introduce others: cost,
  delay, and the fact that a voter cannot verify their own ballot once it enters
  the box. Establish that the problem is real in both directions.}

\subsection{Why a blockchain is the right tool here}

\stub{Argue this honestly. A blockchain is the right tool when multiple mutually
  distrusting parties need a shared, append-only record that no single one of them
  controls. Say where this is genuinely true for voting, and be explicit about
  what a blockchain does \emph{not} fix: voter eligibility still depends on an
  off-chain identity check, which remains a trusted step.}

% ===========================================================================
\section{Requirement Gathering}
% Source material: docs/architecture.md sections 1 and 2
% The template grades "actors/stakeholders' requirement gathering with
% justification", so every requirement below needs a stated reason.
% ===========================================================================

\subsection{Actors and stakeholders}

Table~\ref{tab:actors} lists each actor, what the system permits them to do, and
what it structurally prevents them from doing. The final column is the important
one: a permission model is only meaningful when it is stated in terms of what a
privileged party \emph{cannot} do.

\begin{table}[htbp]
  \centering
  \caption{Actors, their capabilities, and their structural limits.}
  \label{tab:actors}
  \small
  \begin{tabular}{@{}L{2.8cm}L{6.0cm}L{6.4cm}@{}}
    \toprule
    \textbf{Actor} & \textbf{Can do}                                                                                  & \textbf{Cannot do} \\
    \midrule
    Election Authority
                   & Create an election, set candidates and time windows, publish metadata, trigger finalization
                   & Cast or alter a ballot; change the voter root once registration closes; stop the count                                \\
    Registrar
                   & Verify eligibility off-chain, collect identity commitments, publish the roll's root and leaf set
                   & See how anyone voted; add voters after the root is frozen                                                             \\
    Voter
                   & Register an identity commitment, cast exactly one ballot, verify that their ballot landed
                   & Vote twice; vote without being in the committed set                                                                   \\
    Candidate
                   & Be listed on a ballot, publish a manifesto
                   & Anything privileged                                                                                                   \\
    Auditor / Observer
                   & Re-derive the full tally from chain data, verify every proof, check the roll size
                   & Alter anything                                                                                                        \\
    Relayer
                   & Pay gas on a voter's behalf
                   & Forge, modify, or replay a ballot                                                                                     \\
    \bottomrule
  \end{tabular}
\end{table}

\stub{For each actor, write a short paragraph justifying \emph{why} they need
  the capabilities listed and why the restrictions are the right ones. This is the
  part the rubric asks to be justified, so do not just restate the table.}

\subsection{Functional requirements}

\stub{Number these (FR1, FR2, \dots) so Section 5 and the Evaluation 2 report can
  refer back to them. Derive each from an actor in Table~\ref{tab:actors} rather
  than inventing them, e.g. FR1: the system shall allow the Election Authority
  to define a candidate set and a voting window before registration opens.}

\subsection{Non-functional requirements}

\stub{Number these too (NFR1, NFR2, \dots). Cover at least: ballot secrecy,
  eligibility enforcement, individual and universal verifiability, cost per vote,
  latency, and the requirement that a voter need hold no cryptocurrency.}

\subsection{Trust model and explicit non-goals}

A design is only credible if it states what it does \emph{not} achieve.

\stub{Write up the guarantees: eligibility, one-voter-one-ballot, ballot secrecy,
  individual verifiability, universal verifiability, integrity of the count.}

\stub{Write up the non-guarantees honestly, with the reason each is out of scope:
  coercion resistance; Sybil resistance at the registrar; privacy from the relayer,
  which sees the ballot payload it forwards; network-level anonymity; and secrecy
  of the running tally while voting is open.}

% ===========================================================================
\section{Proposed System Design}
% Source material: docs/architecture.md sections 3-7, docs/tech-stack.md
% ===========================================================================

\subsection{Architecture overview}

\stub{Redraw the component diagram from \texttt{docs/architecture.md} as a proper
  figure and place it here. Walk the reader through one request path rather than
  describing every box in isolation.}

\begin{figure}[htbp]
  \centering
  \fbox{\parbox[c][5.5cm][c]{0.9\linewidth}{\centering\itshape
      Figure placeholder: system architecture.\\
      Export to \texttt{figures/architecture.pdf} and replace this box with\\
      \texttt{\textbackslash includegraphics[width=0.9\textbackslash linewidth]\{architecture\}}}}
  \caption{\projectname{} component architecture.}
  \label{fig:architecture}
\end{figure}

\stub{Explain the single most important structural property: the public audit
  path reads the chain directly and never depends on the project's own backend, so
  the result remains verifiable even if the operators are dishonest or offline.}

\subsection{Smart contract design}

\stub{Describe each contract and its responsibility: the election lifecycle state
  machine, the factory that deploys elections as minimal proxy clones, the voter
  registry that stores roots, the tally library, and the ERC-2771 forwarder that
  lets voters transact without holding gas.}

\stub{Draw the election lifecycle as a state machine figure: Created, then
  Registration, Voting, Tallying, Finalized, with Cancelled reachable only before
  voting opens. State which transitions are time-gated and which require an
  explicit call.}

\subsection{Anonymous eligibility and double-vote prevention}

\stub{Explain the mechanism in plain language before any notation. A voter holds
  a secret; its commitment sits as a leaf in a Merkle tree whose root is on-chain.
  To vote, the voter proves in zero knowledge that they know a secret matching
  \emph{some} leaf, without revealing which. The proof carries a nullifier derived
  deterministically from the secret and the election identifier, so a second ballot
  from the same voter is rejected while the chain still learns nothing about who
  they are.}

\subsection{Scalability strategy}

The system is designed so that gas cost is constant in the size of the
electorate. This depends on one decision: the membership tree is built off-chain
and only its root is written on-chain.

\stub{Explain the consequence: a roll of ten voters and a roll of a hundred
  thousand both cost a single 32-byte storage write. State the tradeoff too: the
  tree is not reconstructible from chain data alone, so the registrar must publish
  the leaf set for auditors, and that publication is what keeps the roll auditable.}

\subsection{Technology choices and justification}

\stub{Fill Table~\ref{tab:stack} out into prose. The rubric grades justification,
  so for each row state the alternative that was considered and the specific reason
  it was not chosen.}

\begin{table}[htbp]
  \centering
  \caption{Principal technology choices.}
  \label{tab:stack}
  \small
  \begin{tabular}{@{}L{3.9cm}L{3.9cm}L{7.4cm}@{}}
    \toprule
    \textbf{Layer}     & \textbf{Choice}      & \textbf{Principal reason}                               \\
    \midrule
    Contract language  & Solidity             & \note{ecosystem, EVM portability}                       \\
    Contract toolchain & Foundry              & \note{Solidity-native tests; gas snapshots come free}   \\
    Contract libraries & OpenZeppelin         & \note{audited access control, forwarder, Merkle proofs} \\
    Anonymity          & Semaphore            & \note{audited membership proofs and nullifiers}         \\
    Testbed            & Polygon Amoy         & \note{Sepolia reaches end-of-life 30 Sep 2026}          \\
    Frontend           & Next.js, wagmi, viem & \note{typed contract bindings}                          \\
    Backend            & Fastify, PostgreSQL  & \note{indexer, registrar, relayer}                      \\
    Metadata storage   & IPFS                 & \note{on-chain bytes are expensive}                     \\
    \bottomrule
  \end{tabular}
\end{table}

\stub{Justify the testbed choice explicitly; it is the decision most likely to
  be questioned. Ethereum Sepolia reaches end-of-life on 30 September 2026, which
  falls between this report and the final demonstration, so the project targets
  Polygon Amoy instead and keeps Sepolia only as a comparison baseline.}

% ===========================================================================
\section{Evaluation Parameters}
% Source material: docs/evaluation-parameters.md
% The rubric explicitly grades "defined evaluation parameters", so this section
% carries real weight. Every parameter needs a measurement method, not just a
% name.
% ===========================================================================

Each parameter below states what is measured, how it will be measured, and what
result would count as success. The targets are hypotheses to be tested in
Evaluation 2, not claims; results will be reported as measured, including where
a target is missed.

\subsection{Cost}

\begin{table}[htbp]
  \centering
  \caption{Cost parameters.}
  \label{tab:eval-cost}
  \small
  \begin{tabular}{@{}lL{5.3cm}L{5.1cm}L{3.9cm}@{}}
    \toprule
    \textbf{ID} & \textbf{Parameter}        & \textbf{Measurement}                         & \textbf{Target}       \\
    \midrule
    A1          & Gas per operation         & Gas snapshots over each contract entry point & Ballot under 400k gas \\
    A2          & Cost per vote in INR      & A1 times gas price times token price         & Under INR 0.50        \\
    A3          & Layer-1 vs Layer-2 cost   & Same suite on both networks                  & Quantify the saving   \\
    A4          & On-chain vs IPFS metadata & Inline bytes vs a 32-byte identifier         & Show the gap          \\
    A5          & Proxy deployment saving   & Direct deploy vs minimal proxy clone         & Over 90\% saving      \\
    \bottomrule
  \end{tabular}
\end{table}

\subsection{Performance}

\begin{table}[htbp]
  \centering
  \caption{Performance parameters.}
  \label{tab:eval-perf}
  \small
  \begin{tabular}{@{}lL{5.3cm}L{5.1cm}L{3.9cm}@{}}
    \toprule
    \textbf{ID} & \textbf{Parameter}           & \textbf{Measurement}                   & \textbf{Target}                \\
    \midrule
    B1          & End-to-end vote latency      & Wall clock, click to confirmation      & 95th percentile under 10 s     \\
    B2          & Client proof generation time & Timed in-browser, by device class      & Under 5 s on a mid-range phone \\
    B3          & On-chain verification cost   & Gas attributable to proof verification & Isolate from base cost         \\
    B4          & Throughput                   & Concurrent ballots, ramped             & Locate the knee of the curve   \\
    B5          & Indexer lag                  & Event emission to dashboard visibility & Under 3 s                      \\
    \bottomrule
  \end{tabular}
\end{table}

\subsection{Scalability}

\begin{table}[htbp]
  \centering
  \caption{Scalability parameters.}
  \label{tab:eval-scale}
  \small
  \begin{tabular}{@{}lL{5.3cm}L{5.1cm}L{3.9cm}@{}}
    \toprule
    \textbf{ID} & \textbf{Parameter}       & \textbf{Measurement}          & \textbf{Target}         \\
    \midrule
    C1          & Gas vs voter-roll size   & Rolls from $10^1$ to $10^5$   & Flat                    \\
    C2          & Proof time vs tree depth & B2 across the same roll sizes & Logarithmic, not linear \\
    C3          & Concurrent-voter ceiling & Load until failure            & Name the bottleneck     \\
    C4          & Candidate-count scaling  & 2, 10, and 50 candidates      & Constant per ballot     \\
    \bottomrule
  \end{tabular}
\end{table}

\stub{Say plainly that C1 is the central claim of the design. If gas is not flat
  in the size of the electorate, the architecture did not deliver what it promised,
  and the report will say so.}

\subsection{Security and correctness}

\begin{table}[htbp]
  \centering
  \caption{Security and correctness parameters.}
  \label{tab:eval-sec}
  \small
  \begin{tabular}{@{}lL{5.1cm}L{9.7cm}@{}}
    \toprule
    \textbf{ID} & \textbf{Property}        & \textbf{How it is demonstrated}                                   \\
    \midrule
    D1          & Double-vote resistance   & Replaying a nullifier must revert; fuzzed over collisions         \\
    D2          & Eligibility enforcement  & A valid proof against the wrong root must revert                  \\
    D3          & Ballot secrecy           & Show that no chain data links an address to a ballot              \\
    D4          & Individual verifiability & A voter locates their own ballot transaction                      \\
    D5          & Universal verifiability  & Recompute the tally with the project's backend switched off       \\
    D6          & Authority containment    & The authority cannot vote, alter a frozen root, or stop the count \\
    D7          & Static analysis          & Findings by severity, each with a written disposition             \\
    D8          & Test coverage            & Line coverage over the contract sources                           \\
    D9          & Reorganization safety    & Simulate a chain reorganization; the indexer must self-heal       \\
    \bottomrule
  \end{tabular}
\end{table}

\subsection{Usability and adoption}

\begin{table}[htbp]
  \centering
  \caption{Usability parameters.}
  \label{tab:eval-use}
  \small
  \begin{tabular}{@{}lL{5.3cm}L{5.1cm}L{3.9cm}@{}}
    \toprule
    \textbf{ID} & \textbf{Parameter}     & \textbf{Measurement}                             & \textbf{Target} \\
    \midrule
    E1          & Steps to cast a ballot & User actions from landing page to receipt        & At most 4       \\
    E2          & Voter gas burden       & Cryptocurrency a voter must hold                 & Zero            \\
    E3          & Task success rate      & 10 users unfamiliar with cryptocurrency, unaided & Over 80\%       \\
    E4          & Time to first vote     & Wall clock including wallet installation         & Under 5 minutes \\
    \bottomrule
  \end{tabular}
\end{table}

\stub{Justify why usability is treated as a first-class evaluation parameter
  rather than a nicety: a voting system only redistributes power if ordinary people
  can actually operate it.}

\subsection{Measurement methodology}

\stub{State how measurements will be made credible: repetitions per data point,
  reporting medians and tail percentiles rather than single runs, which environment
  each number comes from, and which controls are recorded alongside each result.}

% ===========================================================================
\section{Implementation Plan}
% Source material: docs/roadmap.md
% ===========================================================================

The system will be built in three functional tiers, each a working system in its
own right. Building them in sequence means every privacy and cost feature is
measured against a working baseline rather than asserted.

\begin{table}[htbp]
  \centering
  \caption{Build tiers.}
  \label{tab:tiers}
  \small
  \begin{tabular}{@{}lL{3.9cm}L{5.3cm}L{3.0cm}@{}}
    \toprule
    \textbf{Tier} & \textbf{Eligibility} & \textbf{Ballot privacy}           & \textbf{Voter pays gas} \\
    \midrule
    1             & On-chain roll        & None; ballots linked to addresses & Yes                     \\
    2             & Merkle root          & None                              & Yes                     \\
    3             & Zero-knowledge proof & Anonymous, nullifier-protected    & No                      \\
    \bottomrule
  \end{tabular}
\end{table}

\begin{table}[htbp]
  \centering
  \caption{Phase plan and milestones.}
  \label{tab:phases}
  \small
  \begin{tabular}{@{}lL{3.0cm}L{10.8cm}@{}}
    \toprule
    \textbf{Phase} & \textbf{Window}       & \textbf{Milestone}                                                        \\
    \midrule
    1              & Late Sep -- mid Oct   & A vote works end to end locally: tier 2, no anonymity, no gas sponsorship \\
    2              & Mid -- late Oct       & Anonymous and gasless on the testbed: tier 3                              \\
    3              & Late Oct -- early Nov & Feature freeze; auditing, resilience, usability testing                   \\
    4              & Early -- mid Nov      & Measurement campaign against Section 5                                    \\
    5              & Mid -- late Nov       & Rehearsal and buffer                                                      \\
    \bottomrule
  \end{tabular}
\end{table}

\stub{Expand each phase into a short paragraph naming its concrete deliverables
  and its exit criterion.}

\stub{State the principal risk and its mitigation: in-browser zero-knowledge
  proof generation is the component most likely to fail, so a dated decision point
  partway through Phase 2 switches to server-side proof generation: weaker
  privacy, honestly documented, rather than allowing the schedule to slip.}

% ===========================================================================
\section{Conclusion}
% ===========================================================================

\stub{Three short paragraphs: what the problem is, what is being built and why
  that design addresses it, and what Evaluation 2 will demonstrate. Do not
  introduce anything new here.}

% ===========================================================================
\section{References}
% ===========================================================================

\printbibliography[heading=none]

\stub{Cite what was actually read. Include the sources behind the design claims
  in Sections 2 and 4, and the documentation for the protocols and testbed named in
  Section 4.}

\end{document}
